\documentclass[a4paper,11pt,oneside]{book}

\usepackage[english]{babel}
\usepackage[usenames,dvipsnames]{xcolor}
\usepackage[utf8]{inputenc}
\usepackage[T1]{fontenc}
\usepackage{changepage}
\usepackage[]{algorithm2e}
\usepackage{amssymb}
\usepackage{amsmath}
\usepackage{graphicx}
\usepackage{listings}
\usepackage{verbatimbox}
\usepackage{setspace} 
\usepackage{ulem}
\usepackage[a4paper,margin=2.5cm,footskip=0.25in]{geometry}
\usepackage{fancyvrb}
\usepackage{float}
\usepackage{hyperref}
\usepackage[parfill]{parskip}
\usepackage{tikz}
\usepackage{pdflscape}
\usepackage{minted}
\usepackage{titlesec}
\usepackage{titleps}
\usepackage{lastpage}
\usepackage{fancyhdr}
\usepackage{etoolbox}
\usepackage{etoolbox}
\makeatletter
\patchcmd{\chapter}{}{}{}{}
\makeatother
\renewcommand{\cleardoublepage}{}
\renewcommand{\clearpage}{}
\renewcommand{\newpage}{}
\onehalfspacing

%New page style for chapters
\newpagestyle{ruledChapter}{
	\setfoot{}{\thepage\ of \pageref{LastPage}}{}
	\footrule
	\renewcommand\makefootrule{\color{black}\rule[\baselineskip]{\linewidth}{0.4pt}}
}
%New page style for rest
\newpagestyle{ruled}{
	\sethead{\raggedright \chaptername\ \thechapter :\ \chaptertitle}{}{}
	\headrule
	\setfoot{}{\thepage\ of \pageref{LastPage}}{}
	\footrule
	\renewcommand\makeheadrule{\color{black}\rule[-.3\baselineskip]{\linewidth}{0.4pt}}
	\renewcommand\makefootrule{\color{black}\rule[\baselineskip]{\linewidth}{0.4pt}}
}

\expandafter\def\csname PY@tok@err\endcsname{}
\newcommand{\HRule}{\rule{\linewidth}{0.5mm}}
\newcommand{\specialcell}[2][c]{%
  \begin{tabular}[#1]{@{}c@{}}#2\end{tabular}}

\addtocontents{toc}{\protect\thispagestyle{empty}}

\title{}
\author{}
\date{} 

\begin{document}
\begin{titlepage}
\begin{center}

%-----------------------------------------------------------------
%							FRONTPAGE
%-----------------------------------------------------------------
\thispagestyle{empty}
\includegraphics[width=0.55\textwidth]{logo.pdf}\\[1cm]    
\textsc{\Large DM848 Microservice}\\[0.5cm]

% Title
\begin{Huge}
\textbf{WebApplication using the microservice architecture}
\end{Huge}

\vspace{1cm}

% Author and supervisor
\begin{minipage}{1\textwidth}
\begin{center}


Lars \textsc{Thomasen}\\
\verb!<latho12@student.sdu.dk>!\\

\end{center}
\end{minipage}


\renewcommand{\contentsname}{Table of Contents}
\tableofcontents
\end{center}
\end{titlepage}


%-----------------------------------------------------------------
%						  ACTUAL REPORT
%-----------------------------------------------------------------
\pagestyle{ruled}

%
% INTRODUCTION
% Contains all nessecary information needed in order to understand the 
% remaining report.
%

\chapter{Introduction}
\setcounter{section}{1}

For this project a web application using the microservice architectural style has been created. This allows one to take the theories learned from the course and implement them into a real product, thus ensuring that the theory is well understood and that the know-how on how to use them in practice is there.

The microservice architectural style is a style to building a single application, consisting of a suite og small services, each running on its own process and instead communicates with the other services. This is a great contrast to the otherwise monolithic style that are well known, where the entire application often consists of a single executable.

This approach offers great flexibility, making it possible to quickly increase the processing power of pre-defined areas of the application, by simply increasing the number of that specific service. Likewise it increases the reliability, as the decline of one or more services wont result in loss of service, as long as atleast one or more services remain running.

\section{Problem formulation}
The specific problem formulation for this project are as following:

\textbf{Goal:} Create a web application using the microservice architecture, this application must have several different services, each serving its own purpose. The application must be highly scalable and take advantage of load balancing and have a well defined interface to the public.

\textbf{Tools:} In order to follow the ideas used widely in the community, this will be done using the Netflix Open Source\footnote{See \textbf{Common Runtime Services \& Libraries} at https://netflix.github.io/} stack, which is a set of tools developed by Netflix when creating microservice services.

\section{Running the code}
The following dependencies are required in order to run the program:
\begin{enumerate}
\item Java 8
\item Maven
\end{enumerate}

A batch script has been created, simply run \texttt{mvnpackage.sh} in order to compile all the services using maven, and then run \texttt{runServices.sh} in order to start the 5 services. Note that starting all the services and registering them with eureka may take a couple of minutes. The webapplication is available at localhost:9999.

A readme file is included in the project root folder if a manual approach is wanted. This allows for setting a given port by feeding it a parameter, and will have logging enabled. Alternatively, the application is deployed at \textbf{\texttt{http://lthomasen.me:9999}} running on an Digital Ocean Ubuntu server.

%
%MODELLING
%
\chapter{Modelling}
This chapter will describe the design choices followed when modelling the application.

\section{Exposing the front-end}
One of the very first thing to define was how to expose all the services to the user. There er several approaches to this, some of these are described here, and the reason for going with one of these are also outlined.

\begin{enumerate}
\item The first and most simple approach would be to have a single monolithic style front-end, which handles all incoming and outgoing requests. This front-end will handle exposing the views and then communicate with all the services behind it. This allows for a static front-end which will always be available at a known location, much like any other approach.

\item Another method would be to split up the front-end into individual services, the best way to describe this would be to imagine the front-end web design split up into sections, and then each section being its own service.

\item And finally, the third option would be to use a proxy, which takes incoming requests and then delegates these requests, often in a round-robin style, to all the services which accepts the type of request given.
\end{enumerate}

For this project option (3) was used. There are several advantages to this approach, and some of these are as following.

The first advantage being that having an extremely lightweight proxy delegate the requests allows for load balancing the web-server, this allows to spin up more than a single web-server such that no bottleneck would be present in generating views and sending them to the users.

The second advantage of this approach is that following the idea of splitting up the front-end into several services would still be possible, by modifying the proxy to simple look at the request type and send it to the correct service.

And finally this allows for having a single point of entry, such that all requests are served over the same port. This avoids any issues which would otherwise be present when working with Ajax-requests using JavaScript, which would be blocking any cross-site scripting. 


\section{The services}
Microservice allows a large system to build up from a number of collaborating components, called services. When creating monolithic structures, we already follow design principles such as Seperation of Concerns (SoC), such that each section of the program addresses a separate well-defined concern.

In micro-services we simply take this idea one step further, and take this entire section and moves it to its own application instead. A big advantage to this is that it gives the option of running additional services of the sections that are expected to be heavily loaded. Thus it is important to also identify which sections we should separate into their own services. Separate too big a section, and we will waste resources spinning up additional services, which are heavier on the resources than they needed to be.

In this application, five services has been created:
\begin{enumerate}
\item \textbf{Registration service}: A service which purpose serves as a central repository, which allows other services to query it such that another service can see that they exists.
\item \textbf{Edge service}: This is the proxy service, which takes incomming requests and redirects them to the correct service. If more than one exists of a given type of service, it will cycle between them in a round-robin manner, such that the load should be evenly divided.
\item \textbf{User service}: Service which handles all operations regarding the users in the system.
\item \textbf{Image service}: Service which handles all operations regarding the images in the system.
\item \textbf{Comment service}: Service which handles all operations regarding the comments in the system.
\end{enumerate}

%
%IMPLEMENTAION
%
\chapter{Implementation}
The implementation of the project is done using the Spring framework, along with the libraries from the Netflix Open Source stack. Spring offers a great framework which simplifies the process of web-applications greatly, by removing a lot of the boilerplate code, and instead attempts to auto-bind elements based on the naming conventions. Spring has little effect on the actual micro-service architecture, but it does allow an easy integration with the Netflix libraries which they have integrated.

Much of the relevant and interesting implementation process for this course were thus understanding how these libraries work, and how we could use them to work together in order to reach the goal set.

\section{Service registration}
Netflix offers a library called Eureka\footnote{A quick glance of what eureka is can be found at https://github.com/Netflix/eureka/wiki}. Eureka is essentially a microservice in itself, its purpose is for locating services and load-balancing. When a service starts, it will attempt to contact the Eureka service and register itself, this is done dynamically at runtime, thus no changes are needed to be done to this service if the load-out of the other services change.

When a service is registered at Eureka, we can simply ask this service whenever we wish to use a specific service, and it will give us the correct address to the service we want. This way the services can have dynamically addresses, and no code will have to be changed if service change location. The only exception here is that the Eureka service itself has a static address for the other services being able to register at it (this can be avoided using a proxy).

Each service has the address of the eureka server configured in the \texttt{*.yml} config files located under the resources folder.  Here we specify the address to the eureka server and the name and port of the offered service. The library then handles the actual registration process, thus no more code will be needed.

\begin{minted}[frame=lines,linenos=true,fontsize=\footnotesize,]{c}
eureka:
  client:
    serviceUrl:
      defaultZone: http://localhost:9998/eureka/
  instance:
    metadataMap:
      instanceId: ${spring.application.name}:
      		${spring.application.instance_id:${server.port}}$
\end{minted}
\section{The proxy}
In order to allow a single point of entry, and making it possible to have several web services available, a proxy has been used. Eureka offers a library known as Zuul, which can work closely together with Eureka. 

\section{Communication}

%
%CONCLUSION
%
\chapter{Conclusion}


%-----------------------------------------------------------------
%						     APPENDIX
%-----------------------------------------------------------------
%\newpage
%\newgeometry{left=2.5cm,right=2.5cm}
%\chapter{Appendix}
%\section{Source code}

\end{document}
